% Baccalauréat malgache, série D, 2010 : corrigé.
\TCDBacAttribution{D}{2010}
\providecommand{\Plan}{(\mathscr{P})}
\providecommand{\Ouv}{(O ; \vec{u} , \vec{v})}
\providecommand{\Courbe}{(\mathscr{C})}
\providecommand{\cm}{~\mathrm{cm}}
\providecommand{\dx}{~\mathrm{d}x}
\providecommand{\Arg}[1]{\mathrm{Arg}\left( #1 \right)}
%
%%%%%%%%%%%%%%%%%%%%%%%%%%%%%%%%%%%%%%%%%%%%%%%%%%%%%%%%%%%%%%%%%%%%%%%%%%%%%%%%%%%%%%%%%%%%%%%%%%%%%%%%%%%%%%%%%%%%%%%%%%%%%%%%
\noindent \underline{\textsc{\textbf{Exercice I}}} (5 points)
\begin{enumerate}
   \item \vln{Image du triangle $SRD$ :}
      La similitude $\sigma$ de centre $R(1 - i)$, d'angle $-\frac{\pi}{2}$ et de rapport $\frac{1}{2}$ s'\'ecrit
      $z' - (1 - i) = -\frac{i}{2}\big(z - (1 - i)\big)$.
      \begin{itemize}
         \item $\sigma(R) = R$ (centre) ;
         \item $\sigma(S)$ : $-\frac{i}{2}(-5 + i) + 1 - i = \frac{1}{2} + \frac{5}{2}i + 1 - i = \frac{3}{2} + \frac{3}{2}i$ ;
         \item $\sigma(D)$ : $-\frac{i}{2}(-1 + 4i) + 1 - i = 2 + \frac{i}{2} + 1 - i = 3 - \frac{1}{2}i$.
      \end{itemize}
      \par\noindent\ansbox{L'image de $SRD$ est le triangle $S'RD'$ avec $S'\left(\frac{3}{2} ; \frac{3}{2}\right)$, $R(1 ; -1)$, $D'\left(3 ; -\frac{1}{2}\right)$.}
      (Figure apr\`es la question 2.)
   \item
      \begin{enumerate}
         \item \vln{Interpr\'etation de $\psi$ :}
            $\psi = \dfrac{z - z_D}{z - z_S}$ (pour $M \neq S$). Donc $|\psi| = \dfrac{MD}{MS}$ et, pour $M \neq S$ et $M \neq D$,
            $\arg\psi = \left(\overrightarrow{MS} ; \overrightarrow{MD}\right)$ $[2\pi]$.
         \item \vln{Ensemble $(\Delta)$ :}
            $|\psi| = 1 \iff MD = MS$ (avec $M \neq S$). \par\noindent\ansbox{$(\Delta)$ est la m\'ediatrice de $[SD]$.}
            En coordonn\'ees : $(x + 4)^2 + y^2 = x^2 + (y - 3)^2 \iff 8x + 6y + 7 = 0$.

            \vln{Ensemble $(\Sigma)$ :}
            $\Arg{\psi} = \frac{\pi}{2}$ signifie $M \neq S$, $M \neq D$ et $\left(\overrightarrow{MS} ; \overrightarrow{MD}\right) = \frac{\pi}{2}$ $[2\pi]$ :
            le triangle $MSD$ est rectangle en $M$, donc $M$ est sur le cercle de diam\`etre $[SD]$ (centre $\Omega\left(-2 ; \frac{3}{2}\right)$,
            rayon $\frac{5}{2}$), et l'orientation de l'angle s\'electionne l'un des deux demi-cercles.
            \par\noindent\ansbox{$(\Sigma)$ est un demi-cercle ouvert de diam\`etre $[SD]$ (priv\'e de $S$ et $D$).}

            \emph{Lequel ?} Pour $z = 0$ : $\psi = \frac{-3i}{4}$, d'argument $-\frac{\pi}{2}$. L'origine est bien sur le cercle de diam\`etre $[SD]$
            (car $\overrightarrow{OS} \perp \overrightarrow{OD}$), mais sur l'autre demi-cercle : \ansbox{$(\Sigma)$ ne passe pas par $O$.}
            $(\Sigma)$ est donc le demi-cercle situ\'e de l'autre c\^ot\'e de la droite $(SD)$ par rapport \`a $O$ ; par exemple, pour
            $z = -\frac{7}{2} + \frac{7}{2}i$, on trouve $\psi = i$.
         \item \vln{Construction :}
            \begin{center}
            \begin{tikzpicture}[scale=0.8,>=stealth]
               \draw[very thin,gray!30] (-5,-2) grid (4,5);
               \draw[->] (-5,0) -- (4,0) node[below left] {$x$};
               \draw[->] (0,-2) -- (0,5) node[below left] {$y$};
               \foreach \t in {-4,-3,-2,-1,1,2,3} \draw (\t,0.06) -- (\t,-0.06) node[below] {\tiny $\t$};
               \foreach \t in {-1,1,2,3,4} \draw (0.06,\t) -- (-0.06,\t) node[left] {\tiny $\t$};
               \begin{scope}
                  \clip (-5,-2) rectangle (4,5);
                  \draw[loko4,thick] plot[domain=-5:4] (\x,{(-7-8*\x)/6});
                  \draw[gray,dashed] (-2,1.5) circle (2.5);
                  \draw[loko1,very thick] (-2,1.5) ++(36.87:2.5) arc (36.87:216.87:2.5);
               \end{scope}
               \draw[loko3,thick] (-4,0) -- (1,-1) -- (0,3) -- cycle;
               \draw[loko6,thick,fill=loko6!15] (1.5,1.5) -- (1,-1) -- (3,-0.5) -- cycle;
               \draw[loko1,fill=white] (-4,0) circle (2pt); \draw[loko1,fill=white] (0,3) circle (2pt);
               \node[below left] at (-4,0) {$S$}; \node[above right] at (0,3) {$D$};
               \node[below left] at (1,-1) {$R$};
               \node[above] at (1.5,1.5) {$S'$}; \node[right] at (3,-0.5) {$D'$};
               \fill (-2,1.5) circle (1.5pt) node[below] {\scriptsize $\Omega$};
               \node[loko1] at (-4.3,3.3) {$(\Sigma)$};
               \node[loko4] at (-4.4,4.4) {$(\Delta)$};
            \end{tikzpicture}
            \end{center}
            (Le triangle $S'RD'$ est l'image de $SRD$ par la similitude de la question 1. L'\'enonc\'e demande l'unit\'e $2\cm$ ; la figure est r\'eduite.)
      \end{enumerate}
   \item
      \begin{enumerate}
         \item \vln{Calcul de $Q(3i)$ :}
            $(3i)^4 = 81$, $-3i(3i)^3 = -3i(-27i) = -81$, $8i(3i) = -24$. Donc $Q(3i) = 81 - 81 - 24 + 24$, \ansbox{$Q(3i) = 0$.}
         \item \vln{Solutions sous forme trigonom\'etrique :}
            On factorise : $(z - 3i)(z^3 + 8i) = z^4 + 8iz - 3iz^3 - 3i\cdot 8i = z^4 - 3iz^3 + 8iz + 24 = Q(z)$.
            $Q(z) = 0 \iff z = 3i$ ou $z^3 = -8i = 8\left(\cos\left(-\frac{\pi}{2}\right) + i\sin\left(-\frac{\pi}{2}\right)\right)$.
            Les racines cubiques de $-8i$ ont pour module $2$ et pour argument $-\frac{\pi}{6} + \frac{2k\pi}{3}$, $k \in \{0, 1, 2\}$,
            soit $-\frac{\pi}{6}$, $\frac{\pi}{2}$, $\frac{7\pi}{6}$ (c'est-\`a-dire $-\frac{5\pi}{6}$).
            \par\noindent\ansbox{$z_0 = 3\left(\cos\frac{\pi}{2} + i\sin\frac{\pi}{2}\right)$, $z_1 = 2\left(\cos\frac{\pi}{2} + i\sin\frac{\pi}{2}\right)$,\\
            $z_2 = 2\left(\cos\left(-\frac{\pi}{6}\right) + i\sin\left(-\frac{\pi}{6}\right)\right)$,
            $z_3 = 2\left(\cos\left(-\frac{5\pi}{6}\right) + i\sin\left(-\frac{5\pi}{6}\right)\right)$.}
            (Soit $3i$, $2i$, $\sqrt{3} - i$ et $-\sqrt{3} - i$.)
      \end{enumerate}
\end{enumerate}
%%%%%%%%%%%%%%%%%%%%%%%%%%%%%%%%%%%%%%%%%%%%%%%%%%%%%%%%%%%%%%%%%%%%%%%%%%%%%%%%%%%%%%%%%%%%%%%%%%%%%%%%%%%%%%%%%%%%%%%%%%%%%%%%
\noindent \underline{\textsc{\textbf{Exercice II}}} (5 points)
\begin{enumerate}
   \item
      \begin{enumerate}
         \item \vln{Fonction de r\'epartition :} $F(x) = P(X \leq x)$ (la somme des probabilit\'es vaut bien $\frac{16}{16}$).
            \[
               F(x) = \begin{cases}
                  0 & \text{si } x < -1, \\[2pt]
                  \dfrac{1}{16} & \text{si } -1 \leq x < 2, \\[4pt]
                  \dfrac{6}{16} = \dfrac{3}{8} & \text{si } 2 \leq x < 3, \\[4pt]
                  \dfrac{9}{16} & \text{si } 3 \leq x < 4, \\[4pt]
                  1 & \text{si } x \geq 4.
               \end{cases}
            \]
         \item \vln{Repr\'esentation graphique} (figure r\'eduite de moiti\'e : $1\cm$ et $4\cm$ au lieu de $2\cm$ et $8\cm$) :
            \begin{center}
            \begin{tikzpicture}[x=1cm,y=4cm,>=stealth]
               \draw[->] (-3,0) -- (6,0) node[below left] {$x$};
               \draw[->] (0,0) -- (0,1.15) node[below left] {$F(x)$};
               \foreach \t in {-1,2,3,4} \draw (\t,0.015) -- (\t,-0.015) node[below] {\scriptsize $\t$};
               \foreach \v/\l in {0.0625/\frac{1}{16},0.375/\frac{3}{8},0.5625/\frac{9}{16},1/1}
                  \draw[gray,dashed] (0,\v) -- (-0.1,\v) node[left] {\scriptsize $\l$};
               \draw[loko1,very thick] (-3,0) -- (-1,0);
               \draw[loko1,very thick] (-1,0.0625) -- (2,0.0625);
               \draw[loko1,very thick] (2,0.375) -- (3,0.375);
               \draw[loko1,very thick] (3,0.5625) -- (4,0.5625);
               \draw[loko1,very thick,->] (4,1) -- (5.8,1);
               \foreach \a/\v in {-1/0.0625,2/0.375,3/0.5625,4/1} \fill[loko1] (\a,\v) circle (2pt);
               \foreach \a/\v in {-1/0,2/0.0625,3/0.375,4/0.5625} \draw[loko1,fill=white] (\a,\v) circle (2pt);
            \end{tikzpicture}
            \end{center}
         \item \vln{Esp\'erance et variance :}
            $E(X) = \frac{-1 + 10 + 9 + 28}{16} = \frac{46}{16}$, \ansbox{$E(X) = \frac{23}{8}$.}
            $E(X^2) = \frac{1 + 20 + 27 + 112}{16} = \frac{160}{16} = 10$, donc
            $V(X) = 10 - \left(\frac{23}{8}\right)^2 = \frac{640 - 529}{64}$, \ansbox{$V(X) = \frac{111}{64}$.}
      \end{enumerate}
   \item Notons $p_J$, $p_G$, $p_D$ les probabilit\'es de prendre deux jaunes, deux gris, deux poussins de couleurs diff\'erentes.
      Ces trois \'ev\`enements sont incompatibles et couvrent tous les cas : $p_J + p_G + p_D = 1$.
      \begin{enumerate}
         \item \vln{$p_G = 0{,}3$ :}
            Les donn\'ees s'\'ecrivent $p_G = 3p_J$ et $p_D = 2p_G = 6p_J$. Donc $p_J + 3p_J + 6p_J = 10p_J = 1$, $p_J = 0{,}1$ et
            \ansbox{$p_G = 0{,}3$.}
            (C'est le cas, par exemple, d'une soubique de 2 jaunes et 3 gris : $p_G = \frac{\mathrm{C}_3^2}{\mathrm{C}_5^2} = \frac{3}{10}$.)
         \item \vln{Au moins un gris :}
            C'est le contraire de << deux jaunes >> : $1 - p_J = 1 - 0{,}1$, \ansbox{la probabilit\'e vaut $0{,}9$.}
      \end{enumerate}
\end{enumerate}
%%%%%%%%%%%%%%%%%%%%%%%%%%%%%%%%%%%%%%%%%%%%%%%%%%%%%%%%%%%%%%%%%%%%%%%%%%%%%%%%%%%%%%%%%%%%%%%%%%%%%%%%%%%%%%%%%%%%%%%%%%%%%%%%
\noindent \underline{\textsc{\textbf{Probl\`eme}}}
\begin{enumerate}
   \item \vln{Ensemble de d\'efinition :}
      Pour $0 < x < 1$, $x > 0$ et $1 - x > 0$ : $\ln x$ et $\ln(1 - x)$ existent ; $f(0)$ et $f(1)$ sont donn\'es.
      \par\noindent\ansbox{$D_f = [0 ; 1]$.}
   \item \vln{Continuit\'e en $0$ et en $1$ :}
      En $0^+$ : $x\ln x \to 0$ (croissances compar\'ees) et $(1 - x)\ln(1 - x) \to 1 \times \ln 1 = 0$, donc $f(x) \to 0 = f(0)$.
      En $1^-$ : on pose $t = 1 - x \to 0^+$ ; $f(x) = -(1 - t)\ln(1 - t) - t\ln t \to 0 = f(1)$.
      \par\noindent\ansbox{$f$ est continue en $0$ et en $1$.}
   \item \vln{Taux d'accroissement :}
      Pour $0 < x < 1$ : $\dfrac{f(x) - f(0)}{x} = -\ln x - \dfrac{(1 - x)\ln(1 - x)}{x}$.
      Or $\dfrac{\ln(1 - x)}{x} = -\dfrac{\ln(1 - x)}{-x} \to -1$ quand $x \to 0$, donc le second terme tend vers $-(1)(-1) = 1$,
      tandis que $-\ln x \to +\infty$ : \ansbox{$\lim\limits_{x\to0^+}\dfrac{f(x) - f(0)}{x - 0} = +\infty$.}

      Remarquons que $f(1 - x) = f(x)$ pour tout $x \in [0 ; 1]$ (la courbe est sym\'etrique par rapport \`a la droite $x = \frac{1}{2}$).
      Avec $t = 1 - x$ : $\dfrac{f(x) - f(1)}{x - 1} = \dfrac{f(t)}{-t} = -\dfrac{f(t) - f(0)}{t - 0} \to -\infty$ quand $t \to 0^+$ :
      \ansbox{$\lim\limits_{x\to1^-}\dfrac{f(x) - f(1)}{x - 1} = -\infty$.}

      \vln{Cons\'equences :}
      \ansbox{$f$ n'est d\'erivable ni en $0$ (\`a droite) ni en $1$ (\`a gauche).}
      La courbe admet aux points $O(0 ; 0)$ et $P(1 ; 0)$ des \emph{demi-tangentes verticales} (port\'ees par les droites $x = 0$ et $x = 1$).
   \item
      \begin{enumerate}
         \item \vln{D\'eriv\'ee :}
            $f'(x) = -\ln x - x\cdot\frac{1}{x} + \ln(1 - x) + (1 - x)\cdot\frac{1}{1 - x} = -\ln x - 1 + \ln(1 - x) + 1$,
            \par\noindent\ansbox{$f'(x) = \ln(1 - x) - \ln x$ pour $0 < x < 1$.}
         \item \vln{In\'equation :}
            Sur $]0 ; 1[$, $\ln$ \'etant strictement croissante : $\ln(1 - x) \leq \ln x \iff 1 - x \leq x \iff x \geq \frac{1}{2}$.
            \par\noindent\ansbox{$S = \left[\frac{1}{2} ; 1\right[$.}
         \item \vln{Tableau de variation :}
            $f' \geq 0$ sur $\left]0 ; \frac{1}{2}\right]$ (nul seulement en $\frac{1}{2}$) et $f' \leq 0$ sur $\left[\frac{1}{2} ; 1\right[$.
            $f\left(\frac{1}{2}\right) = -\frac{1}{2}\ln\frac{1}{2} - \frac{1}{2}\ln\frac{1}{2} = \ln 2 \simeq 0{,}7$.
            \[
               \begin{array}{c|ccccc}
                  x & 0 & & \frac{1}{2} & & 1 \\ \hline
                  f'(x) & \| \; +\infty & + & 0 & - & -\infty \; \| \\ \hline
                  & & & \ln 2 & & \\
                  f & & \nearrow & & \searrow & \\
                  & 0 & & & & 0
               \end{array}
            \]
            (Les symboles $+\infty$ et $-\infty$ dans la ligne de $f'$ indiquent les limites des taux d'accroissement aux bornes.)
      \end{enumerate}
   \item \vln{Trac\'e} (figure r\'eduite ; unit\'e de l'\'enonc\'e : $10\cm$) :
      Tangente horizontale $y = \ln 2$ au point $\left(\frac{1}{2} ; \ln 2\right)$ ; demi-tangentes verticales en $O$ et en $P$.
      \begin{center}
      \begin{tikzpicture}[scale=6,>=stealth]
         \draw[very thin,gray!30,step=0.1] (-0.05,-0.05) grid (1.1,0.85);
         \draw[->] (-0.05,0) -- (1.1,0) node[below left] {$x$};
         \draw[->] (0,-0.05) -- (0,0.85) node[below left] {$y$};
         \foreach \t in {0.5,1} \draw (\t,0.01) -- (\t,-0.01) node[below] {\scriptsize $\t$};
         \draw (0.01,0.6931) -- (-0.01,0.6931) node[left] {\scriptsize $\ln 2$};
         \draw[loko4,thick] (0.3,0.6931) -- (0.7,0.6931);
         \draw[loko4,thick] (0,0) -- (0,0.25);
         \draw[loko4,thick] (1,0) -- (1,0.25);
         \draw[loko1,very thick,smooth,samples=200] plot[domain=0.0005:0.9995] (\x,{-\x*ln(\x)-(1-\x)*ln(1-\x)});
         \fill (1,0) circle (0.25pt) node[below right] {$P$};
         \node[loko1] at (0.83,0.6) {$\Courbe$};
      \end{tikzpicture}
      \end{center}
   \item
      \begin{enumerate}
         \item \vln{D\'eriv\'ees de $u$ et $v$ :}
            $u'(x) = x\ln x + \frac{x^2}{2}\cdot\frac{1}{x} - \frac{x}{2}$, \ansbox{$u'(x) = x\ln x$.}

            $v'(x) = (x - 1)\ln(1 - x) + \frac{(x - 1)^2}{2}\cdot\frac{-1}{1 - x} - \frac{2x - 2}{4}
                   = (x - 1)\ln(1 - x) + \frac{x - 1}{2} - \frac{x - 1}{2}$ (car $\frac{(x-1)^2}{1-x} = 1 - x$),
            \par\noindent\ansbox{$v'(x) = (x - 1)\ln(1 - x) = -(1 - x)\ln(1 - x)$.}
         \item \vln{Primitive cherch\'ee :}
            $f(x) = -x\ln x - (1 - x)\ln(1 - x) = -u'(x) + v'(x)$ : les primitives de $f$ sur $]0 ; 1[$ sont les fonctions
            $F = -u + v + C$, $C \in \mathbb{R}$.
            $u\left(\frac{1}{2}\right) = -\frac{\ln 2}{8} - \frac{1}{16}$ et $v\left(\frac{1}{2}\right) = -\frac{\ln 2}{8} + \frac{3}{16}$, donc
            $-u\left(\frac{1}{2}\right) + v\left(\frac{1}{2}\right) = \frac{1}{4}$. La condition $F\left(\frac{1}{2}\right) = -\frac{3}{4}$ donne $C = -1$.
            Ainsi $F(x) = -\frac{x^2}{2}\ln x + \frac{x^2}{4} + \frac{(x - 1)^2}{2}\ln(1 - x) - \frac{x^2 - 2x}{4} - 1$, soit, apr\`es simplification,
            \[
               \boxed{F(x) = -\frac{x^2}{2}\ln x + \frac{(1 - x)^2}{2}\ln(1 - x) + \frac{x}{2} - 1.}
            \]
      \end{enumerate}
\end{enumerate}
